\cotefile{80}
\scanpage{3}
\mspage{1}

\diagram{\begin{tikzcd}[ampersand replacement=\&,column sep=small]
\deleted{\substack{C(I)\\ \simeq C_{\mathrm{aff}}(I_0)}}\arrow[r]\arrow[d] \& \deleted{C_{n+1}\simeq C(I)/\mathfrak S_I}\\
\deleted{C_{\mathrm{aff}}(n)\simeq C_{\mathrm{aff}}(I_0)/\mathfrak S_{I_0}\simeq C_{\mathrm{aff}}(I)/\mathfrak S_I} \&
\end{tikzcd}}

\deleted{$I,\ i_0\in I$\quad $\operatorname{Card}I=n+1$\quad $I_0=I-\{i_0\}$}

\margin{Classification pseudo-droites}

\diagram{\resizebox{\linewidth}{!}{\begin{tikzcd}[ampersand replacement=\&,column sep=small,row sep=large]
\substack{C(I)\\ \simeq C_{\mathrm{aff}}(I_0)}
 \arrow[dr,"\substack{\text{oubli du }I\text{-}\text{épinglage}\\ \text{sauf }\Delta_{i_0}=\Delta}","\alpha\,\mathord{/\mkern-5mu/}"'] \& \&\\
\& \substack{C_{\mathrm{aff}}(n)\simeq C_{\mathrm{aff}}(I_0)/\mathfrak S_{I_0}\\
 \simeq C_{\deleted{\unsure{\mathrm{aff}}}}(I)/\mathfrak S_{I_0}}
 \arrow[dl,"\beta"'," \substack{=\ \text{oubli du }\Delta\ \unsure{\text{marquage}|\text{marqué}}\\ \deleted{\unsure{\text{gommage du}|\text{quotienté du}}}\\\text{par }\Delta}"]
 \arrow[dr,"\beta'"',"\substack{=\deleted{\text{oubli du }\Delta}\\\text{gommage}\ \text{du }\Delta}" description] \&\\
 C(n+1)\simeq C(I)/\mathfrak S_I \& \& C(n)
\end{tikzcd}}}

\margin{classes d'iso de systèmes $(X,\Sigma,\Delta)$ où $(X,\Sigma)$ syst. de ps dr et $\Delta\in\Sigma$}

$\beta$ surjective, la fibre en \deleted{\unsure{$(X,$|$X$}} $\mathrm{cl}_{n+1}(X,\Sigma)$ est l'ens. des types d'éléments de $\Sigma$, i.e.

\[
\Sigma/G_{X,\Sigma}
\]

[Si $n=5$, $n+1=6$, il y a deux fibres qui sont de card $1$, deux de card. $2$

$\beta'$ surjective (?) La fibre de $\mathrm{cl}_{n}(X,\Sigma^{*})$ est l'ens des types de positions relatives p.r. à $(X,\Sigma^{*})$

[Si $n=5$, $C(n)$ a un élément, la fibre est de card. \unsure{$6$|$5$}]

\uline{NB}\quad $C_{\mathrm{aff\,st.}}(I_0)\subset C_{\mathrm{aff}}(I_0)$.\quad stable par $\mathfrak S_I$, a fortiori par $\mathfrak S_{I_0}$, passant aux quotients on trouve les diagrammes analogues avec

\[
C_{\mathrm{st}}(I)\quad\text{''}\qquad C\text{''}(I)
\]

\[
C_{\mathrm{st}}(I)=C_{\mathrm{aff\,st}}(I_0),\ C_{\mathrm{aff\,st}}(n),\ C_{\mathrm{st}}(n),\ C_{\mathrm{st}}(n+1).
\]
% archive sheet number: 2
